\documentclass[11pt,a4paper]{article}
\usepackage[margin=2.2cm]{geometry}
\usepackage[T1]{fontenc}
\usepackage{amsmath,amssymb,xcolor,enumitem,booktabs,graphicx,tabularx,array}
\usepackage[normalem]{ulem}
\setlength{\parskip}{0.6em}\setlength{\parindent}{0pt}
\definecolor{newc}{HTML}{1f6f43}
\newcommand{\new}[1]{\textcolor{newc}{#1}}
\newcommand{\old}[1]{\textcolor{gray}{\sout{#1}}}
\newcommand{\note}[1]{{\small\itshape [#1]}}
\newcommand{\lab}[1]{\textsf{\footnotesize\detokenize{#1}}}
\renewcommand{\ref}[1]{\lab{#1}}
\renewcommand{\eqref}[1]{(\lab{#1})}
\newcommand{\nref}[1]{\lab{#1}}
\newcommand{\citep}[1]{[\lab{#1}]}
\newcommand{\citet}[1]{\lab{#1}}
\newcommand{\where}[1]{\par\textbf{#1}\par}
\title{Proposed change, Problem Formulation, round 2\\ What else moves out of the chapter}
\date{25 November 2026, version 1}
\begin{document}\sloppy
\maketitle

\section*{How to read this}
[A133] is in the thesis. This round takes the ten further points you chose,
1 to 10. Each one either cuts a design reason that the Method already gives,
moves it there, or moves reporting detail to the Experiments or Appendix C.
Black text stays, \new{green} is new and \old{struck} text goes. \note{Notes
in brackets} say why. Answer by number. Nothing below is in the thesis yet.

\begin{tabular}{@{}llp{0.6\linewidth}@{}}
\toprule
Item & To & What \\
\midrule
R1 & Method (already there) & how the moves are numbered \\
R2 & Method & why messages are learned \\
R3 & Method & why the field of view is small \\
R4 & Method (already there) & why one policy is shared \\
R5 & Method & what $\theta$ stands for per controller \\
R6 & Method & why every controller reads the crowding feature \\
R7 & Method & smoothing of $\hat\omega$ \\
R8 & Experiments & stalls \\
R9 & Experiments & the backlog shown over the run \\
R10 & Appendix C & how the keep-up check counts, with the two count examples \\
\bottomrule
\end{tabular}

The chapter shrinks by about a page. The Method grows by about a third of a
page, since most of what goes there is already in it.

% ================================================================
\section*{Where the items are}

\subsection*{R7. Smoothing of $\hat\omega$}
\where{Problem Formulation, ``What the rule reads''}
Demand storage reads only the first number. Congestion-aware storage reads
both. \old{A shelf sees only a few picks in one storage period, so a single
period gives a number that jumps up and down by chance. $\hat\omega_t(v)$
therefore looks back over several of the shelf's recent periods. A shelf
with no pick yet reads 0.} Formally, the order estimate is \dots

\where{Method, ``Where arriving stock goes'', the paragraph before \lab{eq:standstill}}
$\hat\omega_t(v)$ is the plain average of the shelf's last five readings.
Five is a practical smoothing choice. \new{A shelf sees only a few picks in
one period, so one reading jumps up and down by chance.} It evens out one
noisy period and still follows recent congestion, and the number was not
tuned. A shelf with no pick in a period keeps its average\new{, and a shelf
with no pick yet reads 0}.

\note{What $\hat\omega$ means stays in the Problem Formulation. How it is
smoothed is a choice the Method already makes and defends.}

% ================================================================
\section*{Who decides}

\subsection*{R4. Why one policy is shared}
\where{Problem Formulation, ``Decentralised''}
A decentralised controller decides for one agent, from what that agent sees
around it, its \emph{observation} $o_i(t)$. Every agent uses the same
policy. When communication is switched on, an agent also gets learned
messages from the agents near it. \old{Sharing one policy suits agents that
are all alike. It does not make the decisions central, because each agent
still acts only on what it sees and what nearby agents send it.} \new{The
shared policy does not make the decisions central, because each agent still
acts only on what it sees and what nearby agents send it.} Each agent
chooses on its own, \dots

\note{Why a policy is shared is a training choice. ``How the controllers
learn'' gives it, citing the Background. Nothing is added to the Method.}

% ================================================================
\section*{What a decentralised agent sees}

\subsection*{R3. Why the field of view is small}
\where{Problem Formulation, after \lab{fig:observation}}
\old{A small field of view keeps the policy's input the same size in any
warehouse, so a policy trained on one warehouse can run on a bigger one.
But the}\new{The} agent cannot see a conflict coming from further away than
$d_{\mathrm{obs}}$.

\where{Method, ``How an agent chooses its move'', second paragraph}
The network follows \citet{vanKnippenberg2021}, who built on PRIMAL
\citep{sartoretti2019primal}. It keeps their local view and their distance
label. \new{A local view keeps the input the same size in any warehouse, so
a policy trained on one warehouse can run on a bigger one.} It replaces
their graph convolution \dots

\note{The limit stays in the problem. The reason for a small view is a
network design reason.}

\subsection*{R6. Why every controller reads the crowding feature}
\where{Problem Formulation, after the crowding equation}
\old{A controller that can tell how crowded its area is can route around the
crowd before an agent gets stuck in it, the same argument
\citet{pham2024optimizingcrowdawaremultiagentpath} make for their
crowd-density feature. Here every controller gets the number, not only the
decentralised one, so it is measured the same way for every agent. Every
controller reads only the agent's own value. A zone or the centralised
controller sees more agents, so their values could be worked out, but no
controller reads them. Appendix~\ref{app:controllers} explains why a
section-based agent's field of view is not cut at the zone edge.}
\new{Every controller reads the crowding feature, not only the decentralised
one, and each reads only the agent's own value.}
Figure~\ref{fig:crowding} gives each agent's reading on the small graph.

\where{Method, ``What the network reads''}
The crowding feature $\delta_i(t)$ is added once for the agent, not for each
vertex. \new{With it a controller can route around a crowd before an agent
gets stuck in it, the argument
\citet{pham2024optimizingcrowdawaremultiagentpath} make for their
crowd-density feature. Every controller gets it, so it is measured the same
way for every agent. A zone or the centralised controller sees more agents,
so their values could be worked out, but no controller reads them.
Appendix~\ref{app:controllers} explains why a section-based agent's field of
view is not cut at the zone edge.} The vertices chosen by earlier agents
\dots

\note{The problem keeps the fact that every controller reads it, which the
heading also says. The why and the source go to where the network reads
it.}

\subsection*{R2. Why messages are learned}
\where{Problem Formulation, ``Messages from nearby agents''}
\old{The content of a message is learned, not chosen by hand. Each sender
computes its message from its own observation with the same network that
chooses its action, and training decides what is worth sending, as in
\citet{pham2024optimizingcrowdawaremultiagentpath} and \citet{li2022pico}. A
hand-picked message, such as the sender's distance to its target, cannot say
which way the sender is heading. Because a learned message depends on the
network, it is part of the policy and not of the observation, so $o_i(t)$
keeps its three pieces.} \new{What a message says is learned, so it is part
of the policy and not of the observation, and $o_i(t)$ keeps its three
pieces.} Whether messages help at all is the question RQ4
(Section~\ref{sec:pd-rq}). With communication switched off, the agent acts
on $o_i(t)$ alone.

\where{Method, ``Messages'', after the first sentence}
\textbf{Messages.} A message has no layer of its own, and this is the
thesis's own choice. \new{What a message says is learned, not chosen by
hand. A hand-picked message, such as the sender's distance to its target,
cannot say which way the sender is heading.} In
\citet{pham2024optimizingcrowdawaremultiagentpath} an agent's state comes
from an image network \dots

\note{Who can hear whom (the communication graph) and why only the
decentralised controller has messages stay, because those are limits of the
problem. That messages are trained through the policy loss, as in Pham et
al.\ and Li et al., the Method already says at the end of the same
paragraph.}

% ================================================================
\section*{The learning problem}

\subsection*{R1. How the moves are numbered}
\where{Problem Formulation}
A stochastic game gives every agent one fixed list of actions, but the
legal moves depend on where the agent stands. So the list is as long as the
busiest vertex needs, one wait and one move per outgoing edge\new{, and
moves that lead nowhere are switched off}. \old{Move number $k$ takes the
$k$-th outgoing edge of the vertex the agent stands on, in vertex-number
order. Moves that lead nowhere are switched off before the policy chooses.
If the busiest junction has four exits, every agent has five actions, and in
a one-way aisle two of them are switched on.} Formally, with $d_{\max}$ the
largest number of outgoing edges at any vertex, the action list is
$\mathcal{U}_i = \{\mathrm{wait},\, \mathrm{move}_1, \dots,
\mathrm{move}_{d_{\max}}\}$.

\where{Method, ``Only moves that exist''}
With $d_{\max} = \max_{v\in V} |N^{+}(v)|$ there are $d_{\max}+1$ slots.
\new{If the busiest junction has four exits, every agent has five slots,
and in a one-way aisle two of them are open.} At a vertex with fewer edges,
the empty slots get $-\infty$ before the softmax \dots

\note{The numbering in vertex order is already in the Method (``one for each
outgoing edge in increasing vertex number''). Only the example moves.}

\subsection*{R5. What $\theta$ stands for}
\where{Problem Formulation, after the objective}
\old{The same objective is used for all three learned controllers. Only what
$\theta$ stands for changes. For the decentralised controller it is the
weights shared by all agents. For the centralised and section-based
controllers it is the weights of the network over the whole graph or over a
zone, as in \nref{eq:centralisedpolicy} and \nref{eq:sectionpolicy}.}

\old{Sharing one objective does not mean the three controllers learn the
same thing.} \new{The objective is the same for the three controllers, but
they do not learn the same thing.} An action is judged by the average reward
of the agents that chose together with it. So a decentralised agent is not
charged for the delay it causes others, but a centralised agent is.

\where{Method, ``How the controllers learn'', after ``share one policy''}
All agents of a controller share one policy, the common recipe of
Section~\ref{sec:bg:marl}. \new{Its weights $\theta$ are the network run by
every agent of the decentralised controller, or the network over the whole
graph or a zone for the other two. Every controller is trained on the same
objective \eqref{eq:objective}.} Each controller is trained separately
\dots

\note{The line above the cut paragraph already says $\pi_\theta$ stands for
$\pi^{\mathrm{cen}}_\theta$, $\pi^{\mathrm{sec}}_\theta$ or
$\pi^{\mathrm{dec}}_\theta$. What the weights are is a training fact. The
difference in credit stays, because it is part of what the comparison
compares.}

% ================================================================
\section*{How a run is scored}

\subsection*{R10. How the keep-up check counts}
\where{Problem Formulation}
The check counts every task the fleet finishes in the measured steps, also
one released before them\old{. A fleet that fell behind earlier still works
through its older tasks, and that work is what shows whether the queue grows
or holds. So}\new{, so} the measure is the share of the arriving work the
fleet serves. \old{With 900 tasks released in the measured steps and 600
tasks finished in them, the share is two thirds and the queue grows by
300.}

A pick counts as released when its order arrives, also while its stock is
not yet on a shelf\new{, because the order is work the fleet owes
(Appendix~\ref{app:measures:keepup})}. \old{The order is work the fleet
owes, and the backlog already counts it. Left out, a fleet whose put-aways
fail would release almost no picks and seem to keep up. Such a fleet with
800 orders waiting for stock and 10 picks finished serves about one pick in
eighty.}

\where{Appendix C, new subsection after the crowd density paragraph of \lab{app:measures}}
\new{\textbf{How the keep-up check counts} \note{label
\lab{app:measures:keepup}}}

\new{The check \eqref{eq:throughput} counts every task the fleet finishes in
the measured steps, also one released before them. A fleet that fell behind
earlier still works through its older tasks, and that work is what shows
whether the queue grows or holds. With 900 tasks released in the measured
steps and 600 tasks finished in them, the share is two thirds and the queue
grows by 300.}

\new{A pick counts as released when its order arrives, also while its stock
is not yet on a shelf. The order is work the fleet owes, and the backlog
already counts it. Left out, a fleet whose put-aways fail would release
almost no picks and seem to keep up. Such a fleet with 800 orders waiting
for stock and 10 picks finished serves about one pick in eighty.}

\note{The chapter keeps both rules in one sentence each and the formal
definition. The reasons and the two count examples move together, so the
appendix reads on its own. The one pointer goes to an appendix, which you
allow.}

\subsection*{R9. The backlog shown over the run}
\where{Problem Formulation}
The check looks at one finite horizon, so passing it does not prove that
the backlog stays bounded. A run can finish most of its tasks while the
backlog still grows slowly. \old{The backlog over the run is therefore shown
for the conditions in the main result, so a growing backlog is visible.}

\where{Experiments, ``Scoring a condition'', after ``reported next to it''}
\new{The check looks at one horizon, so the backlog over the run is shown for
the conditions in the main result, and a backlog that still grows is
visible.}

\note{The limit belongs to the measure and stays. What the Results show is
reporting.}

\subsection*{R8. Stalls}
\where{Problem Formulation}
\old{A fleet can also jam for a stretch and then catch up. Its throughput
recovers, so it passes the check, and its service time carries the jam.
The jam is counted on its own as a stall. A stall is a stretch after the
warm-up in which no agent finishes a task for longer than a set length
(Appendix~\ref{app:cal:stall}). A light load leaves longer quiet stretches,
so the length grows with the time between tasks. It looks at the whole fleet
and not at one agent, because one agent often waits by plan, as under Token
Passing, while a healthy fleet finishes tasks all the time. It is counted
because it says something about the controller rather than about the
load.}

\where{Experiments, new paragraph after ``Scoring a condition''}
\new{\textbf{Stalls.} A fleet can also jam for a stretch and then catch up.
Its throughput recovers, so it passes the check, and its service time
carries the jam. The jam is counted on its own as a stall. A stall is a
stretch after the warm-up in which no agent finishes a task for longer than
a set length (Appendix~\ref{app:cal:stall}). A light load leaves longer
quiet stretches, so the length grows with the time between tasks. It looks
at the whole fleet and not at one agent, because one agent often waits by
plan, as under Token Passing, while a healthy fleet finishes tasks all the
time. It is counted because it says something about the controller rather
than about the load.}

\note{Unchanged text, moved. The stall is not in \lab{tab:measures} and is
not part of the decision problem. It is a count the Results report, set in
Appendix H, so it sits with how a condition is scored.}

% ================================================================
\section*{When approved}
Applied under [A134], with a GAPS.md entry. R10 adds the label
\lab{app:measures:keepup} to Appendix C, so \texttt{tools/sync\_xrefs.py
--build} is run. Every tag goes at the end of a line, and I grep every new
tag line for text after the \texttt{\%} before building.

\end{document}
